\section{Lean formalization outline and path to object iPRA}\label{sec:formalization}

This section specifies the representations, dependency order, induction
invariants, and final theorem types needed for a complete formalization.  The
files in \code{lean/} mirror the plan.  They are an integration scaffold, not a
claim that the analytic development has already been checked by Lean.

\subsection{The endpoint that counts}

At the pinned ProveIt revision, the ordinary iPRA syntax is in
\code{LeanProofs.ConciseIPRA}.  Primitive-recursive symbols have type
\code{PR k}; terms have type \code{PTm n}; the axioms are \code{PRAAx}; and the
convenient induction rule is \code{PRA.qfInd}.  A completed theorem must have
the shape
\begin{lstlisting}[language={}]
import ConciseIPRA.PR
import ConciseIPRA.NativeBiInterpretation

namespace LeanProofs.ConciseIPRA

-- phi : Fm PTm 0
example : [] |- [PRAAx] phi := by
  ... -- an explicit object-language derivation

end LeanProofs.ConciseIPRA
\end{lstlisting}
After this theorem exists, \code{pra\_iff\_native} transfers it to the checked
native scalar calculus.  A theorem merely about Lean's \code{Nat},
\code{Rat}, or a host complex type is useful evidence but is not a proof in
iPRA.

Every formalized result should therefore have three visibly separate layers:
\begin{enumerate}
\item an executable host definition and regression theorem;
\item a primitive-recursive graph formula or bounded certificate checker;
\item a theorem of type \code{[] |- [PRAAx] phi} proving that checker sound.
\end{enumerate}

\subsection{Compilation rather than a universal evaluator}

The user-facing notation permits variable lists of parameters, but iPRA need
not quantify over an all-purpose program interpreter.  A small external
compiler takes a concrete theorem family---fixed $p,q,m,n$, fixed formulas for
the parameter names, and fixed branch/tail algorithms---and emits ordinary
\code{PR} symbols.  List coding, bounded loops, rational arithmetic, and
finite-polynomial operations are themselves primitive recursive.  Soundness
is proved once for their checkers.  The resulting theorem is a schema over
fixed compiled names, exactly as in the existing Gamma and Barnes-$G$ packets.

For a genuinely uniform family such as ``all rational parameter lists of
length at most $r$'', $r$ is external and the generated checker uses bounded
loops through $r$.  This avoids both higher types and a universal
primitive-recursive evaluator.

\subsection{Data types and numerical predicates}

The host scaffold uses the following records.
\begin{lstlisting}[language={}]
structure QComplex where
  reNum : Int
  reDen : Nat       -- positive
  imNum : Int
  imDen : Nat       -- positive

structure MeijerSpec where
  aNum : List QComplex
  aDen : List QComplex
  bNum : List QComplex
  bDen : List QComplex

structure BranchCert where
  centerArg : Rat
  radiusArg : Rat
  awayZero  : Rat
  cutMargin : Rat

structure EvalCert where
  truncation : Nat
  guardBits  : Nat
  tailNum    : Nat
  tailPow    : Nat
  clearance  : Rat
\end{lstlisting}
The names \code{aNum}, \code{aDen}, \code{bNum}, and \code{bDen} denote the
four parameter blocks $A^-,A^+,B^-,B^+$.  In the object language these records
are natural-number tuples obtained by pairing.

A complex ball uses the $\ell^1$ norm because it avoids square roots:
\[
 B(c,p)=\{w:\abs{\Re(w-c)}+\abs{\Im(w-c)}\le2^{-p}\}.
\]
The object predicate \code{BallEq(x,y,p)} is a cleared integer inequality
between the centers and error numerators returned by two fixed evaluators.
There is no set-valued object in the formal language.  Division is never a
primitive theorem conclusion: a reciprocal operation requires
\code{Apart(q,d)}, and all final inequalities are cross-multiplied by the
positive denominator $d$.

The basic reusable object lemmas are:
\begin{enumerate}
\item signed-integer and rational normalization, order, and absolute value;
\item complex addition, multiplication, integer powers, and finite products;
\item ball addition, multiplication, Horner evaluation, and telescoping sums;
\item reciprocal with a quantitative apartness witness;
\item dyadic exponential, logarithm, trigonometric, and square-root names;
\item bounded sums, products, search, pairing, and simultaneous recursion;
\item finite Laurent polynomials with truncated convolution.
\end{enumerate}
The checked \code{IPRADerivations} arithmetic library supplies much of the
first and sixth items.  The Gamma packet is the intended source of the
certified Gamma, reciprocal-Gamma, logarithm, and Stirling interfaces, but its
analytic contracts still have to be instantiated as object derivations rather
than imported as axioms.

\subsection{Module graph}

A dependency order minimizing circularity is
\begin{verbatim}
MeijerG/Codes
  -> MeijerG/RationalBalls
  -> MeijerG/FiniteLists
  -> MeijerG/HypergeometricTerms
  -> MeijerG/HypergeometricTail
  -> MeijerG/GammaRatio
  -> MeijerG/SlaterRight
  -> MeijerG/SlaterLeft
  -> MeijerG/LaurentJets
  -> MeijerG/ResidueEvaluator
  -> MeijerG/ContourEvaluator
  -> MeijerG/ContourResidueAgreement
  -> MeijerG/DifferentialEquation
  -> MeijerG/ParameterAlgebra
  -> MeijerG/Derivatives
  -> MeijerG/Mellin
  -> MeijerG/RationalPowerProducts
  -> MeijerG/Laplace
  -> MeijerG/TransformKernels
  -> MeijerG/Whittaker
  -> MeijerG/SpecialCases
  -> MeijerG/ObjectTargets
  -> MeijerG/Checks/Axioms
\end{verbatim}
The hypergeometric and residue modules should be completed before the contour
modules.  This makes most source-page identities available without waiting for
a general complex-integration library.

\subsection{Stage 1: exact hypergeometric terms}

For externally fixed numerator parameters $\alpha_1,\ldots,\alpha_r$ and
denominator parameters $\beta_1,\ldots,\beta_s$, define a term graph by
\begin{align*}
 T_0&=1,\\
 T_{k+1}\,(k+1)\prod_{j=1}^{s}(\beta_j+k)
 &=T_k z\prod_{i=1}^{r}(\alpha_i+k).
\end{align*}
The object implementation stores a rational complex center and a common
positive denominator.  The primary theorem is the cleared recurrence; a
quotient presentation is derived only under denominator apartness.

The partial-sum loop maintains the quantifier-free invariant
\begin{equation}\label{eq:lean-hyp-invariant}
 S_k=\sum_{j<k}T_j,
 \qquad
 D_kT_{k+1}=N_kT_k,
 \qquad
 \operatorname{Err}_k\le E_k,
\end{equation}
where all equalities are represented by integer cross products.  If a bounded
checker proves $\abs{T_{j+1}}\le\rho\abs{T_j}$ for $j\ge N$ and
$\rho\le1-2^{-h}$, the geometric tail theorem yields
\[
 \abs{\sum_{j\ge N}T_j}\le \frac{\abs{T_N}}{1-\rho}.
\]
The proof is qf induction on the finite geometric sum, followed by a cleared
inequality.  This module proves \cref{thm:hyp-conv} for every compiled series.

\subsection{Stage 2: Gamma ratios and Slater coefficients}

A Slater term contains products such as
\[
 C_h=
 \frac{\prod_{j\le m,j\ne h}\Gamma(b_j-b_h)
       \prod_{j\le n}\Gamma(1+b_h-a_j)}
      {\prod_{j>m}\Gamma(1+b_h-b_j)
       \prod_{j>n}\Gamma(a_j-b_h)}.
\]
At requested precision $p$, the compiler allocates guard bits backwards
through the finite product tree, requests Gamma balls, checks denominator
apartness, and outputs a ball for $C_h$.  The soundness invariant records a
product center, an accumulated numerator error, and a positive lower bound for
every denominator factor.

The right-Slater target is schematically
\begin{lstlisting}[language={}]
RightAdmissible spec z p ->
BallEq
  (meijerRight spec z (p + guard spec p))
  (sumOverRightPoles spec z (p + guard spec p))
  p
\end{lstlisting}
and similarly for the left expansion.  The proof has two parts:
\begin{enumerate}
\item exact residue algebra for a finite closing rectangle;
\item the explicit closing-side and series-tail inequalities from
      \cref{sec:convergence}.
\end{enumerate}
For most special cases only the second line---the already normalized
hypergeometric sum---is needed as the definition, so an early formalization can
prove all coefficient identities before contour agreement.

\subsection{Stage 3: confluent poles as finite Laurent jets}

Encode a truncated Laurent polynomial by a list indexed from $-r$ to $d$.
Provide PR functions for addition, shift, truncated convolution, and coefficient
selection.  For a pole of multiplicity $r$, the invariant after multiplying
$i$ factors is
\[
 F_i(u)=\sum_{j=-r}^{r-1}c_{i,j}u^j+u^rR_i(u),
 \qquad \abs{R_i(u)}\le M_i
 \quad(\abs u\le\rho).
\]
The update for $M_i$ is a finite Cauchy-product bound.  Multiply by the finite
jet
\[
 z^{s_0+u}=z^{s_0}\sum_{j<r}\frac{(u\Log z)^j}{j!}+u^rR_z(u)
\]
and select coefficient $u^{-1}$.  The output is a polynomial in $\Log z$ of
degree at most $r-1$, exactly as in \cref{def:jet}.

The key object theorem is not ``take a limit as parameters collide''.  It is a
finite coefficient theorem saying the emitted jet is the residue of the
certified local product.  This handles integer-order $K_\nu$, $Y_\nu$, the
six-term identity in \cref{sec:elementary-cases}, and every polynomial repair.

\subsection{Stage 4: the contour evaluator}

A contour is a finite list of rational complex vertices plus, for an unbounded
vertical contour, a truncation function $T(p)$.  On each segment, use a
composite Taylor or trapezoidal rule with a supplied derivative bound.  The
object predicate \code{ContourSafe} checks:
\begin{enumerate}
\item positive distance from every relevant pole and from zero;
\item a branch box containing the image of $z^s$;
\item the vertical estimate \eqref{eq:master-vertical-bound};
\item a quadrature bound for every segment;
\item a closing-side bound for the finite residue rectangle.
\end{enumerate}
The finite residue theorem from the ProveIt Cauchy--residue packet can then be
specialized to the Gamma-ratio integrand.  Its output is the
\code{ContourResidueAgreement} theorem.  This theorem is optional for the
coefficient identities but essential for the literal Mellin--Barnes definition
and transform formulas.

\subsection{Stage 5: differential equation without differentiation axioms}

There are two complementary formalizations.

First, for a residue series $F(z)=z^\lambda\sum c_kz^k$, calculate the
coefficient of $z^{\lambda+k}$ in
\[
 (-1)^{p-m-n}z\prod_j(\thetaop-a_j+1)F
 -\prod_j(\thetaop-b_j)F.
\]
It is exactly the cleared hypergeometric term recurrence.  A qf induction proves
that every partial polynomial satisfies the equation up to an explicit tail;
the tail tends below $2^{-p}$ by the same ratio certificate.

Second, on a finite contour, applying $\thetaop$ multiplies the kernel by $s$.
Gamma recurrence converts the operator multiplier to zero identically.  The
finite quadrature error is propagated through the fixed polynomial in $s$.
This proves the ODE on contour domains and agrees with the series theorem on
overlap.

A representative object target is
\begin{lstlisting}[language={}]
Admissible spec z p ->
BallEq
  (meijerODELeft spec z (p + guard))
  (meijerODERight spec z (p + guard))
  p
\end{lstlisting}
where both sides are finite combinations of derivative-name evaluators; no
quantification over an arbitrary differentiable function occurs.

\subsection{Stage 6: parameter algebra and derivative rules}

Cancellation, power shift, block permutations, inversion, and integer block
moves are finite identities of Gamma products.  Represent a parameter list as
a bounded array and normalize it to a multiset of tagged affine expressions.
The certificate consists of a permutation, a cancellation matching, and a
small exponent adjustment.  Its checker proves equality of the kernels by
bounded product commutativity.

For the multiplication theorem, Gauss multiplication expands each Gamma
factor into $k$ factors.  The certificate records the affine map
$a\mapsto(a+r)/k$ or $b\mapsto(b+r)/k$, the extracted power of $k$, and the
argument scale.  All exponents are rational and are checked after clearing the
common denominator.

Derivative identities use a finite polynomial in $s$ and one Gamma recurrence.
The primary formal statement is a coefficient or contour ball equality, not an
unbounded limit definition of derivative.  A generic fixed-order Taylor
contract then identifies that name with the derivative on a certified disc.

\subsection{Stage 7: Mellin, Euler, rational powers, Laplace, and transform kernels}

A finite Fubini certificate contains four truncation bounds, a common
quadrature mesh, and a rational majorant for the discarded rectangles.  The
proof first rearranges a genuinely finite double sum.  It then adds:
\begin{enumerate}
\item $x$-endpoint tails;
\item contour-frequency tails;
\item quadrature errors;
\item the inner Beta or Gamma integral error.
\end{enumerate}
The Mellin transform theorem uses the triangular Fourier/Fej\'er kernel from
\cref{def:mellin-cert}; its normalization and positivity are separate finite
lemmas.  The final theorem is parameterized by a fixed function name and its
explicit modulus.  No theorem quantifies over ``all integrable functions''.

The Laplace theorem is easier: the inner integral is a Gamma contract and the
remaining Gamma product is normalized syntactically to the target $G$.
Convolution concatenates two parameter arrays after applying the Mellin theorem
at $1-s$.  These normalized-list equalities should be checked by an executable
certificate, avoiding large manual rewrites in Lean.

For rational powers $x^\gamma,x^\delta$, normalize
$\gamma/\delta=A/B$, substitute $s=Bu$, and expand every slope-$A$ or
slope-$B$ Gamma factor with the finite Gauss multiplication theorem.  The
certificate stores the expanded parameter blocks, the collected rational
power, and the prefactor.  Its soundness theorem is a bounded-list induction
plus the existing Gamma-multiplication contracts.

Transform-pair modules never quantify over an unspecified function space.  A
fixed test-function code carries endpoint, Mellin, Fubini, and inversion
moduli; the checker verifies that the two finite Mellin multipliers multiply to
one.  The Whittaker module separately compiles the $U$ and ${}_1F_1$
reductions used by the forward and inverse kernels, including branch and Gamma
denominator guards.

\subsection{Stage 8: special functions}

Each elementary or higher special case has a very small certificate:
\begin{enumerate}
\item a concrete \code{MeijerSpec};
\item its single- or double-residue expansion;
\item a list of Gamma cancellations;
\item a coefficient recurrence matching the target series;
\item a branch equality such as $\sqrt{x^2}=x$;
\item a common tail bound.
\end{enumerate}
For example, the cosine proof compiles to
\begin{lstlisting}[language={}]
SafeCosBranch x p ->
BallEq
  (sqrtPi * meijerG cosSpec (x*x/4) (p+g))
  (cosName x (p+g))
  p
\end{lstlisting}
The induction invariant states equality of the first $N$ rational terms; the
factorial duplication theorem converts the coefficient denominators.  The
Lerch theorem uses the exact identity
\code{poch a k / poch (a+1) k = a / (a+k)} in cross-multiplied form.  Bessel
$Y$ and integer-order $K$ additionally invoke the jet module.

\subsection{Representative object-language target list}

The companion \code{ObjectLanguageTargets.lean} records the following final
families.
\begin{enumerate}
\item \code{HyperTermRecurrence}: exact cleared recurrence for every compiled
      hypergeometric term.
\item \code{HyperTail}: a ratio certificate implies the requested tail ball.
\item \code{SlaterRight} and \code{SlaterLeft}: residue evaluator equals the
      compiled hypergeometric sum.
\item \code{ContourEqualsResidues}: the finite contour evaluator agrees with
      the residue evaluator.
\item \code{MeijerODE}: the two sides of the differential equation agree.
\item \code{PowerShift}, \code{Inversion}, \code{CancelTopBottom}, and
      \code{MultiplyParameters}.
\item Integral targets \code{Mellin}, \code{Euler}, \code{Convolution},
      \code{RationalPower}, and \code{Laplace}, each under an explicit
      certificate predicate.
\item Transform targets \code{KernelPair}, \code{Narain},
      \code{GenLaplace}, and \code{Whittaker} for fixed coded test classes.
\item \code{ExpCase}, \code{CosCase}, \code{LogCase},
      \code{IncompleteGammaCase}, \code{BesselCases}, and \code{LerchCases}.
\end{enumerate}
All equality targets are \code{BallEq} schemas in the requested precision.
Exact finite identities, such as parameter-list normalization, are proved as
ordinary object equalities and then used inside the ball proofs.

\subsection{Milestones and trust audit}

A practical sequence of independently useful milestones is:
\begin{enumerate}
\item checked exact rational/complex balls and hypergeometric recurrence;
\item checked geometric/factorial tail theorem;
\item cosine, exponential, logarithm, and rational special cases entirely in
      \code{Prf PRAAx};
\item generic nonconfluent Slater evaluator;
\item Laurent-jet residues and Bessel $K,Y$;
\item parameter algebra and the differential equation;
\item contour/residue agreement;
\item Mellin, convolution, and Laplace certificates;
\item transform-pair instances.
\end{enumerate}

Each completed module should be audited for \code{sorry}, \code{admit},
\code{axiom}, \code{Classical.choice}, and untrusted native decision procedures.
Use \code{\#print axioms} on host helper theorems, and additionally inspect the
final theorem type to confirm that it is a \code{Prf PRAAx} derivation.  Python
and Wolfram checks in this archive are regression evidence only; they are not
part of the trusted proof base.
